\textbf{Время:} резерв

\textbf{Тезис:} общий график позволяет сопоставить видимые масштабы и зависимости радиусов SMASH и STAR при сохранении ограничений архивного сравнения.

\textbf{На экране:} шапка Au+Au, 3 ГэВ, 0–10\%, LCMS; верхняя легенда — бирюзовые ромбы SMASH, тёмные заполненные и светлые пустые кружки двух зарядов STAR; панели out, side, long слева направо; две оговорки внизу.

\textbf{Текст:} Здесь три модельных радиуса и результаты STAR поставлены на общие оси. В шапке перечислены общие обозначения: Au+Au, три ГэВ, центральность от нуля до десяти процентов, система LCMS.

Сначала прочитаем легенду. Бирюзовые ромбы — мой расчёт SMASH 2025 года. Заполненные тёмные кружки — пары отрицательных пионов STAR, пустые светлые — положительных. Во всех панелях по горизонтали k-T, по вертикали радиус в фемтометрах.

Пройдём слева направо. На панели out модельные центры лежат выше экспериментальных. Для side последовательности расположены ближе друг к другу. Для long модельные центры также выше. Во всех трёх панелях с ростом k-T видна тенденция к уменьшению радиусов.

Теперь прочитаем нижнюю подпись. Центры точек оцифрованы с исходных рисунков, а условия отбора в модели и эксперименте различаются. Поэтому по этому сравнению можно обсуждать видимые масштабы и зависимости; точное количественное согласие оно не устанавливает. У STAR сохранён исходный статус Preliminary.

Такой график помогает сформулировать следующий шаг: повторить сравнение при согласованных отборах и проверенной обработке.

\textbf{Переход:} Соберём всю цепочку: что мы узнали об источнике, начав только с зарегистрированных частиц?

\textbf{Справка:} Балинов (2025), исходный слайд 17; Anna Kraeva for STAR, Universe 10, 188 (2024), рис. 5, PDF-с. 9, DOI 10.3390/universe10040188. Центры маркеров считаны с исходных растров: 12 модельных и 24 экспериментальных, два заряда STAR отдельно. Это приближённая ручная оцифровка, не исходные численные таблицы и не новый фит; отрезки между модельными центрами служат только направляющими для глаза. Погрешности остаются на исходных графиках и не переносились на общие оси; вывод о согласии в их пределах или оценка качества фита по этому рисунку не допускаются.

Модельный отбор η = −1.5…1.5 отличается от лабораторного отбора STAR −2…0; распределения пар по быстроте также различаются. Последний интервал kT не совпадает. В STAR учтено кулоновское взаимодействие, архивная модельная обработка использует плоские волны; заряд именно этих модельных панелей не установлен. Четыре видимые группы STAR не дополнены пятой, хотя в тексте статьи перечислены пять интервалов; графические смещения двух зарядов сохранены. Код, настройки и физическая корректность численного масштаба архивного фита независимо не подтверждены; ошибки записи формул защиты не объявляются только типографическими. Поэтому даже близость отдельных центров не доказывает правильность динамики модели.

Таблица: figures/comparison\_digitized\_2025\_star.csv. Калибровка, пиксельные координаты, источники и ограничения: figures/comparison\_digitized\_2025\_star\_provenance.json.
